% !TeX program = lualatex
% Compile with LuaLaTeX (Overleaf: Menu > Compiler > LuaLaTeX).
% Step-by-step instructions: GUIDE.md in the template repository.
\documentclass[titlepage]{generalreport} % titlepage: title on its own page,
                                         % numbering starts on page 2

% Anything left out (except \title) is simply not printed.

% ---- Title page head --------------------------------------------------------
\logo{assets/logos/placeholder-logo}   % top right; e.g. \logo[height=1cm]{…}
\emblem{assets/logos/placeholder-emblem} % title-page watermark, bottom right;
                                       % options: [width=…, color=black!8,
                                       % right=.333, bottom=.333]
\infoline{\textit{Kurs oder Veranstaltung}, weitere Angaben}
                                       % top left, free text; wraps at half the
                                       % text width, \\ breaks manually
\date{\today}                          % line below; or e.g. \date{10.10.2026}

% ---- Title block -----------------------------------------------------------
\documenttype{BERICHT}                 % label above the title, as typed
\title[Kurztitel]{Titel des Berichts}  % [short title for the running head]
\subtitle{Untertitel des Berichts}
\author[1]{Vorname Nachname}           % one to three authors; [n] = number
\author[1]{Vorname Nachname}           % of the author's affiliation
\affiliation[1]{Institut, Universität, Stadt} % period is added automatically

\addbibresource{references.bib}

% ---- Example content -------------------------------------------------------
% Provides the German dummy text (\blindtext, \Blindtext) in the sections.
% Delete once the sections contain real text.
\usepackage{blindtext}

\begin{document}

\maketitle

\begin{abstract}
Kurze Zusammenfassung: Ziel, Methode, wichtigste Ergebnisse.
\end{abstract}


\tableofcontents
\clearpage

% One file per section; add, remove or reorder the lines as needed.
\section{Einleitung}
% Beispielinhalt: alles unterhalb der \section-Zeile kann gelöscht werden.

\blindtext

Dieser Absatz verweist auf \cref{fig:a}, auf \cref{eq:lindblad} und auf
\cref{tab:werte}, ohne dabei etwas Bestimmtes auszusagen.%
\footnote{Auch diese Fußnote hat keinen Inhalt.}
Inline-Mathematik wie $e^{i\pi} + 1 = 0$ oder
$i\hbar\,\partial_t \lvert\psi\rangle = \hat H \lvert\psi\rangle$ steht
mitten im Satz, und eine Größe mit Einheit wie
$\qty{42.0(5)}{\kilo\gram\metre\per\second\squared}$ ebenfalls.

\blindtext

\section{Theoretische Grundlagen}
% Beispielinhalt: alles unterhalb der \section-Zeile kann gelöscht werden.

\blindtext{} Ein Literaturverweis steht hier nur zur Zierde
\cite{zinth_optik_2018}.

\subsection{Ein Unterabschnitt}

Ein Blindsatz behauptet, dass sich alles nach der Gleichung
\begin{equation}
  i\hbar\,\frac{\partial}{\partial t}\Psi(\mathbf{r},t)
  = \Bigl[-\frac{\hbar^{2}}{2m}\nabla^{2} + V(\mathbf{r},t)\Bigr]
    \Psi(\mathbf{r},t)
  \label{eq:schroedinger}
\end{equation}
richtet, was aber niemand überprüft hat.

% Gleichungen benennt man im Satz; die Nummer bleibt das einzige Label.
Ein weiterer Blindsatz nennt die Lindblad-Gleichung, die für diesen Text
keinerlei Bedeutung hat:
\begin{equation}
  \frac{\mathrm{d}\hat\rho}{\mathrm{d}t}
  = -\frac{i}{\hbar}\bigl[\hat H, \hat\rho\bigr]
    + \sum_{k} \gamma_{k} \Bigl( \hat L_{k} \hat\rho \hat L_{k}^{\dagger}
    - \frac{1}{2} \bigl\{ \hat L_{k}^{\dagger} \hat L_{k}, \hat\rho \bigr\}
    \Bigr) .
  \label{eq:lindblad}
\end{equation}
Im weiteren Text heißt sie einfach Lindblad-Gleichung \cref{eq:lindblad};
\Cref{eq:schroedinger} am Satzanfang setzt das Wort \enquote{Gleichung}
automatisch davor.

\subsection{Noch ein Unterabschnitt}

\blindtext

Mehrzeilige Ausdrücke wie die Dyson-Reihe werden mit \texttt{align}
gesetzt; jede Zeile erhält eine Nummer, \verb|\notag| unterdrückt sie:
\begin{align}
  \hat U(t, t_{0})
  &= \mathcal{T} \exp\Bigl( -\frac{i}{\hbar} \int_{t_{0}}^{t}
       \hat H(t')\, \mathrm{d}t' \Bigr)
     \label{eq:dyson} \\
  &= \mathbb{1} - \frac{i}{\hbar} \int_{t_{0}}^{t} \mathrm{d}t_{1}\,
       \hat H(t_{1}) \notag \\
  &\quad + \Bigl(-\frac{i}{\hbar}\Bigr)^{2} \int_{t_{0}}^{t} \mathrm{d}t_{1}
       \int_{t_{0}}^{t_{1}} \mathrm{d}t_{2}\, \hat H(t_{1}) \hat H(t_{2})
       + \dots
     \label{eq:dyson-reihe}
\end{align}
Ob \cref{eq:dyson} und \cref{eq:dyson-reihe} hier etwas zu suchen haben,
bleibt offen.

\section{Versuchsaufbau und Durchführung}
% Beispielinhalt: alles unterhalb der \section-Zeile kann gelöscht werden.

\blindtext

\begin{figure}[htbp]
  \centering
  \includegraphics[width=0.6\textwidth]{example-image-a}
  \caption{Eine Abbildung mit dem Buchstaben A, die nichts darstellt}
  \label{fig:a}
\end{figure}

Die folgenden Schritte beziehen sich lose auf \cref{fig:a}:
\begin{enumerate}
  \item Ein erster Schritt ohne Bedeutung.
  \item Ein zweiter Schritt, der dem ersten folgt.
  \item Ein dritter Schritt mit Unterpunkten:
    \begin{enumerate}
      \item ein Unterpunkt,
      \item noch ein Unterpunkt.
    \end{enumerate}
  \item Ein letzter Schritt.
\end{enumerate}

\blindtext

\section{Teilversuch 1}
% Beispielinhalt: alles unterhalb der \section-Zeile kann gelöscht werden.

\blindtext

\begin{figure}[htbp]
  \centering
  \begin{subfigure}{0.45\textwidth}
    \includegraphics[width=\linewidth]{example-image-b}
    \caption{Buchstabe B}
    \label{fig:b}
  \end{subfigure}\hfill
  \begin{subfigure}{0.45\textwidth}
    \includegraphics[width=\linewidth]{example-image-c}
    \caption{Buchstabe C}
    \label{fig:c}
  \end{subfigure}
  \caption{Zwei Teilabbildungen nebeneinander. Lange Beschriftungen werden
    im Blocksatz gesetzt, kurze zentriert}
  \label{fig:bc}
\end{figure}

\Cref{fig:bc} zeigt zwei Buchstaben; \cref{fig:c} lässt sich auch einzeln
referenzieren. Setzt man \cref{eq:dyson} beliebig oft in sich selbst ein,
erhält man angeblich die Werte in \cref{tab:werte}.

\begin{table}[htbp]
  \caption{Zahlen ohne Bedeutung, an der Dezimalstelle ausgerichtet}
  \label{tab:werte}
  \centering
  \begin{tabular}{
      S[table-format=1.2]
      S[table-format=2.2]
      S[table-format=2.1(2)]}
    \toprule
    {$x$ / \unit{\metre}} & {$y$ / \unit{\second}}
      & {$z$ / \unit{\kilo\gram}} \\
    \midrule
    0.25 & 14.47 & 14.1(8)  \\
    0.50 & 28.94 & 29.5(8)  \\
    0.75 & 43.41 & 42.7(9)  \\
    1.00 & 57.88 & 58.6(10) \\
    \bottomrule
  \end{tabular}
\end{table}

\begin{figure}[htbp]
  \centering
  \includegraphics[width=0.7\textwidth]{example-image}
  \caption{Eine breite Platzhalterabbildung}
  \label{fig:d}
\end{figure}

Aus \cref{fig:d} folgt mit etwas Fantasie
$\zeta = \qty{58.1(9)}{\micro\electronvolt\per\tesla}$.

\blindtext

\section{Teilversuch 2}
% Beispielinhalt: alles unterhalb der \section-Zeile kann gelöscht werden.

\blindtext

Der Wert aus \cref{fig:d} weicht um $\qty{0.4}{\percent}$ von einem Wert
ab, der nirgends steht \cite{zinth_optik_2018}. Mögliche Gründe:
\begin{itemize}
  \item ein erster Grund,
  \item ein zweiter Grund mit Unterpunkten,
    \begin{itemize}
      \item ein Unterpunkt,
      \item noch ein Unterpunkt,
    \end{itemize}
  \item ein Grund, der auf \cref{fig:bc} verweist.
\end{itemize}

Das Feynmansche Pfadintegral passt zwar nicht zu \cref{eq:lindblad},
steht aber trotzdem hier:
\begin{equation}
  \langle x_{b}, t_{b} \vert x_{a}, t_{a} \rangle
  = \int \mathcal{D}[x(t)]\,
    \exp\biggl( \frac{i}{\hbar} \int_{t_{a}}^{t_{b}}
    \Bigl[ \frac{m}{2}\dot x^{2} - V(x) \Bigr] \mathrm{d}t \biggr) .
  \label{eq:pfadintegral}
\end{equation}

\blindtext

\section{Zusammenfassung}
% Beispielinhalt: alles unterhalb der \section-Zeile kann gelöscht werden.

\blindtext

Zusammengefasst stehen in \cref{tab:werte} Zahlen, die weder zu
\cref{eq:pfadintegral} noch zu \cref{fig:d} passen; der Code, der sie
erzeugt hat, steht in \cref{sec:notebook}.

\clearpage
\section{Anhang}
% Pro Code-Datei ein Unterabschnitt: ein, zwei Sätze, wofür die Datei da ist,
% dann \codefile[angezeigter Name]{Pfad}{Hinweis}. Zuerst das Notebook,
% danach die Module.
% Notebook exportieren:  jupyter nbconvert --to script auswertung.ipynb
%   (--no-prompt entfernt die "# In[n]:"-Zeilen,
%    --TemplateExporter.exclude_markdown=True die Markdown-Zellen)
% Beispielinhalt: Unterabschnitte und Dateien in code/ können gelöscht werden.

\subsection{Auswertungsnotebook}\label{sec:notebook}
Das Notebook \texttt{auswertung.ipynb} führt die gesamte Auswertung durch und
erzeugt \cref{fig:d}. Markdown-Zellen erscheinen als Kommentare,
\texttt{\# In[n]:} markiert die Zellen.
\codefile[auswertung.ipynb]{code/auswertung.py}{Jupyter-Notebook, exportiert}

\subsection{Modul \texttt{fitting.py}}\label{sec:fitting}
Gewichtete Anpassung einer Ursprungsgeraden; wird in Zelle~3 des
Notebooks (\cref{sec:notebook}) verwendet.
\codefile{code/fitting.py}{Python-Modul}

\subsection{Versuchsprotokoll}\label{sec:protokoll}
Das handschriftliche Versuchsprotokoll ist auf den folgenden Seiten
eingescannt.
% Alle Seiten des Scans (A4), je eine pro Seite in voller Textbreite;
% pagecommand={} behält Kopf- und Fußzeile.
\includepdf[pages=-, width=\textwidth, pagecommand={}]{assets/placeholder-protocol}

\clearpage
\printbibliography[heading=bibnumbered]

\end{document}
